\chapter{Vitamin synthesis and nutrient metabolism}

\section{Learning objectives}
Distinguish endogenous synthesis from activation and recycling; identify the limited vitamin pathways humans possess; explain the role of plants and microorganisms; and understand why minerals are obtained from the environment rather than synthesized by the body.

\section{What ``synthesis'' means in nutrition}
In nutrition, synthesis may refer to several different processes. A cell may build a vitamin molecule from a precursor, convert an ingested precursor into an active form, recycle a cofactor after it has performed its function, or obtain a compound made by microorganisms. These processes are not interchangeable. A precursor pathway may be insufficient under ordinary conditions, and an activated metabolite may have a different regulatory role from the dietary form.

\section{Human synthesis and activation}
Humans cannot synthesize most vitamins in sufficient amounts and therefore require dietary sources. Skin synthesis of vitamin D is a true example of human vitamin synthesis. Other pathways are conversions: the body can make niacin from tryptophan and convert dietary provitamin A carotenoids into vitamin A, but both processes have limits and carotenoid conversion varies. Gut bacteria produce menaquinones (vitamin K2), although how much the body absorbs and uses is uncertain. \cite{vitaminDMetabolism,niacinDRI,odsA,odsK}

Vitamin D provides the clearest example of endogenous synthesis. Ultraviolet B radiation converts 7-dehydrocholesterol in the epidermis to previtamin D3, which rearranges to cholecalciferol (vitamin D3). The liver hydroxylates vitamin D to 25-hydroxyvitamin D, the principal circulating status metabolite. The kidneys and some other tissues can hydroxylate 25-hydroxyvitamin D to 1,25-dihydroxyvitamin D (calcitriol), a hormone-like regulator of calcium and phosphate physiology. Skin pigmentation, age, latitude, season, and the amount of skin exposed influence production. Typical sunscreen use does not appear to eliminate skin production, but ultraviolet exposure carries skin-cancer risk and should not be prescribed casually as a vitamin dose. \cite{vitaminDMetabolism,odsD}

\section{Plant and microbial biosynthesis}
Plants, algae, fungi, and bacteria synthesize vitamins or vitamin precursors that enter human food chains. Plants make carotenoids and folates; yeasts and fungi can make or accumulate B vitamins; and some bacteria synthesize cobalamins or menaquinones. Food processing, fermentation, storage, soil, light, and cooking influence the final amount. Fermentation may improve digestibility or create a bioavailable form, but it does not guarantee that a food supplies a complete nutrient requirement. \cite{odsA,odsFolate,odsB12,odsK}

\section{Vitamin activation and coenzyme formation}
After absorption, vitamins often undergo phosphorylation, methylation, hydroxylation, or conjugation. Thiamin becomes thiamin diphosphate, riboflavin becomes FMN and FAD, niacin becomes NAD and NADP, and B6 becomes pyridoxal phosphate. Folate is converted into tetrahydrofolate derivatives that carry one-carbon units. The active forms of B12 support methionine synthase and methylmalonyl-CoA mutase. These reactions help explain why liver, kidney, intestinal, or mitochondrial dysfunction can alter nutrient metabolism. \cite{odsThiamin,odsRiboflavin,odsNiacin,odsB6,odsFolate,odsB12}

\section{Minerals are not synthesized}
Minerals are chemical elements, not organic molecules, so the human body cannot make them from other substances. It can absorb, transport, store, bind, and excrete them, and some elements can change chemical form or become incorporated into proteins. Plants acquire minerals from soil and water; animals obtain them through food. Geology, agriculture, fortification, water treatment, and food processing therefore influence mineral exposure. \cite{odsIron,napElectrolytes}

\section{Recycling and conservation}
The body conserves some nutrients through recycling and storage. Macrophages recover iron from aging red blood cells; it returns to circulation bound to transferrin or is stored as ferritin. Vitamin K is regenerated during the clotting cycle by vitamin K epoxide reductase, the enzyme inhibited by warfarin. Folate circulates through cellular folate pools, but body stores are limited and can be depleted over months. These processes help conserve nutrients but do not eliminate dietary requirements. Blood loss, inflammation, malabsorption, and medicines can alter nutrient availability or status. \cite{odsIron,odsK,odsFolate}

\section{Industrial production}
Commercial vitamins may be made by chemical synthesis, microbial fermentation, extraction from natural materials, or a combination of these methods. Manufacturing affects a product's chemical form, purity, stability, excipients, and dose accuracy. Vitamin D2 is commonly made by ultraviolet irradiation of yeast ergosterol. Vitamin D3 is typically made by irradiating 7-dehydrocholesterol from sheep's wool (lanolin); lichen-derived, animal-free D3 is also available. Both forms are used in supplements and fortified foods. A label's source alone does not establish superior clinical efficacy. Quality control and contaminant testing matter more than the word ``natural.'' \cite{odsD}

\section{Health effects of impaired synthesis}
Disorders of synthesis, absorption, activation, transport, or regulation can impair vitamin status despite reasonable intake. Examples include reduced skin production of vitamin D, impaired liver or kidney activation, malabsorption, genetic enzyme defects, and effects of some medicines. Clinical management should identify the relevant problem and select treatment for the specific diagnosis; simply increasing a precursor may not correct a downstream defect and can cause toxicity. \cite{odsD,vitaminDMetabolism,odsSafety}

\section{Clinical case}
A housebound older adult has low serum 25-hydroxyvitamin D and chronic kidney disease. The relevant question is not simply whether to take more dietary vitamin D. Assessment should include calcium, phosphate, parathyroid hormone, kidney function, medicines, fracture risk, and the clinician's treatment plan. This illustrates the difference between nutritional intake, endogenous synthesis, activation, and treatment of a metabolic disorder. \cite{odsD,vitaminDMetabolism}

\section{Summary}
Most vitamins must be obtained from food because human synthesis is absent or inadequate, although the body can make or activate a few vitamins from precursors. Minerals cannot be synthesized by the body. Health depends on the entire chain from environmental production and food preparation to intestinal absorption, cellular activation, recycling, and excretion.
